\section{Discussion}\label{sec:discussion}

In this study, we developed a framework for examining how IV assumptions affect American-option values while retaining a chosen physical stock-return model. The paired ablation established that changing the IV module could alter interim marks and liquidation outcomes on identical stock paths, even though entry marks and terminal payoffs agreed exactly. This was a controlled comparison of model behavior. The chronological evaluation asked the separate question of whether the assembled model predicted later observations accurately. The fixed coupled factor produced only small average reductions in mean-based option forecast error relative to frozen IV, and it did not show a consistent advantage across dates or point summaries. Its central intervals also behaved differently across the two tickers. A more elaborate IV process therefore did not automatically provide a better forecast of future option values. The framework's demonstrated contribution was the ability to compare these assumptions while preserving the contracts, stock paths, and pricing calculation. That control made both the measured changes and their limits easier to interpret. It also prevented terminal payoff agreement from being mistaken for validation of interim IV dynamics. The construction remains useful for studying how specified stock and IV assumptions affect simulated option values, but agreement within the simulator does not establish agreement with the market. The observed option errors and stock benchmarks were necessary to assess that additional requirement.

The stock evaluation clarified why this distinction mattered for forward use. Replacing simulated stock paths with the subsequently observed paths reduced option errors on the same contracts and dates, but supplied information unavailable when a forecast would have been issued. This diagnostic identified an important contribution from the stock forecast; it did not offer a usable correction or establish the cause of the stock misses. The subsequent stock-only tests found no return-scaling error, and conditioning the initial state on recent returns did not materially improve five-session central forecasts. The fitted transitions quickly lost information about the initial return state. A larger comparison in 2025 then showed central errors close to an unchanged-price benchmark, while earlier validation preferred zero drift for the tested directional regression. These findings concerned the specified models and inputs, not the general possibility of predicting returns. Adaptive volatility improved LLY's distribution score in both evaluation periods, but did so through a different range of possible prices rather than a learned directional forecast. Its poorer GS scores showed why a wider interval could not be treated as an improvement by itself. The evidence thus supported a limited, ticker-dependent change in uncertainty estimates. It did not justify replacing the stock generator throughout the framework or treating its central projections as reliable predictions of future direction. We retained JumpHMM as the worked example and reported the alternatives as benchmarks.

Separating the stock and IV components also left structural requirements that forecast comparisons could not resolve. A lattice can account for early exercise in an individual contract without ensuring consistent values across strikes. The residual GS violations demonstrated this limitation even under direct evaluation of the fitted surface (Supplementary Table~\ref{tab:strike_audit}). For a strategy involving several contracts, those inconsistencies matter because decisions depend on relative prices as well as the value of each position. Arbitrage-constrained surface methods offer one way to address that requirement before outputs are used as market-consistent training data~\cite{cont2023}. Surface calibration posed another issue. The neural models improved pooled IV fit, yet the May 11 MSFT and LLY panels exposed substantial errors in the date-specific level. The larger July-cutoff fit retained a related problem on later observations. Pooling more dates did not itself supply a mechanism for updating a frozen ticker level. Origin anchoring removed the initial level difference from the forecast comparison, but was an explicit change in initialization and did not impose cross-contract consistency. Repricing with reported endpoint IV further reduced errors for the examined contracts, although it again used later information. Agreement with their quoted spreads did not establish universal compatibility between Alpaca's Black--Scholes IV field and an American lattice. Surface error, stock-distribution accuracy, option-price error, and strike consistency therefore remained distinct requirements. Improving one did not establish the others or supply the joint pricing interpretation of a calibrated stock-volatility model.

The evidence remains limited by the observation windows, contract availability, and scope of model selection. The initial surface comparisons used fifteen snapshots, and the temporal extensions used one fixed neural training seed. The option forecast test covered 23 capture sessions, while the additional stock-only test supplied a full 2025 observation period. The latter did not extend the option validation to 2025. Overlapping horizons and common market movements reduced the independent information in both evaluations; additional simulated paths did not add market observations. Maturity-bucket collection left the originally selected monthly contracts without five-session endpoints. The shorter-maturity cohort was added after an availability audit and did not validate the original month-long illustrations. The stock benchmark classes were proposed after inspecting the 2026 misses, although their settings were selected using earlier data. The tested predictors and penalties were deliberately limited, and no general superiority claim followed from them. The original JumpHMM, IV-factor parameters, and return coupling remained fixed, while the strike audit covered a finite grid and did not establish calendar or dynamic no-arbitrage. Valuation omitted dividends, execution costs, margin, premium financing, hedging, and prior assignment. Retrieval timestamps did not establish exact stock--option synchronization, and state- and mood-dependent IV mechanisms remained inactive. These limits bound the present use of the framework to controlled modeling comparisons and the specific forecasting evidence reported here. Consistent contract histories and further evaluation across market conditions are needed before drawing broader conclusions about practical trading use.
